\chapter{Numerical Methods and Interview Problems}

Direct factorisations (Chapters 4, 11, 17) cost $O(n^3)$ and need the whole matrix in memory. For the huge sparse matrices of scientific computing and machine learning, one instead uses \emph{iterative} methods that only need to multiply a vector by the matrix. This chapter covers the core iterative algorithms: power iteration, the QR algorithm for all eigenvalues, Krylov subspaces and the conjugate gradient method, and randomised SVD. The second half is a collection of worked interview problems that exercise the whole book.

\section{Power iteration and its relatives}

\begin{theorem}[Power iteration]\label{thm:power}
Let $A\in\R^{n\times n}$ be diagonalisable with eigenvalues $\abs{\lambda_1}>\abs{\lambda_2}\ge\dots\ge\abs{\lambda_n}$ and eigenvectors $\mathbf{v}_i$. If $\mathbf{x}_0=\sum c_i\mathbf{v}_i$ has $c_1\ne0$, then the normalised iterates $\mathbf{x}_{k+1}=A\mathbf{x}_k/\norm{A\mathbf{x}_k}$ satisfy
\[ \mathrm{dist}\big(\mathbf{x}_k,\Span\{\mathbf{v}_1\}\big)=O\big(\abs{\lambda_2/\lambda_1}^k\big),\qquad\frac{\mathbf{x}_k\T A\mathbf{x}_k}{\mathbf{x}_k\T\mathbf{x}_k}\to\lambda_1. \]
\end{theorem}

\begin{proof}
$A^k\mathbf{x}_0=\sum c_i\lambda_i^k\mathbf{v}_i=\lambda_1^k\big(c_1\mathbf{v}_1+\sum_{i\ge2}c_i(\lambda_i/\lambda_1)^k\mathbf{v}_i\big)$. The bracket converges to $c_1\mathbf{v}_1$ at the stated rate, and normalising does not change the direction. For symmetric $A$ the Rayleigh quotient is accurate to the square of the vector error.
\end{proof}

Each step costs one matrix--vector product. The method is slow when $\abs{\lambda_2/\lambda_1}\approx1$ and finds only the dominant eigenpair. Two simple modifications repair this.

\begin{itemize}
  \item \textbf{Inverse iteration.} Apply the power method to $(A-\mu I)^{-1}$, whose dominant eigenvalue is $1/(\lambda_j-\mu)$ for the eigenvalue $\lambda_j$ closest to the shift $\mu$. One factorisation of $A-\mu I$ is reused at every step, and convergence is fast when $\mu$ is a good guess.
  \item \textbf{Rayleigh quotient iteration.} Update the shift each step to $\mu_k=\mathbf{x}_k\T A\mathbf{x}_k$. For symmetric matrices this converges cubically.
\end{itemize}

\begin{example}
For $A=\begin{pmatrix}2&1\\1&2\end{pmatrix}$ (eigenvalues $3,1$), starting at $\mathbf{x}_0=(1,0)$: $A\mathbf{x}_0=(2,1)$, then $A(2,1)=(5,4)$, then $(14,13)$, giving directions $(0.89,0.45)$, $(0.78,0.62)$, $(0.73,0.68)\to(0.71,0.71)$, converging at rate $\abs{1/3}^k$ to $(1,1)/\sqrt2$.
\end{example}

\section{The QR algorithm}

To find \emph{all} eigenvalues one uses the QR algorithm, named for the QR factorisation of Chapter~11, and one of the ten most influential algorithms of the twentieth century.

\begin{definition}[Basic QR iteration]
Set $A_0=A$. For $k=0,1,2,\dots$: factor $A_k=Q_kR_k$ and set $A_{k+1}=R_kQ_k$.
\end{definition}

\begin{proposition}[Properties]\label{prop:qralg}
$A_{k+1}=Q_k\T A_kQ_k$, so every $A_k$ is orthogonally similar to $A$ and has the same eigenvalues. Moreover $A_k=(Q_0\cdots Q_{k-1})\T A(Q_0\cdots Q_{k-1})$, and the product $Q_0\cdots Q_{k-1}R_{k-1}\cdots R_0=A^k$ is the QR factorisation of $A^k$.
\end{proposition}

\begin{proof}
$R_k=Q_k\T A_k$, so $A_{k+1}=R_kQ_k=Q_k\T A_kQ_k$. The second statement follows by induction, and for the last: $A^{k+1}=A\cdot A^k=A\,\mathcal{Q}_k\mathcal{R}_k=\mathcal{Q}_kA_k\mathcal{R}_k=\mathcal{Q}_kQ_kR_k\mathcal{R}_k$, where $\mathcal{Q}_k=Q_0\cdots Q_{k-1}$.
\end{proof}

Under mild conditions $A_k$ converges to an (upper) triangular matrix whose diagonal lists the eigenvalues in decreasing order of modulus: the algorithm is the power method on a whole basis at once ($A^k$ applied to an orthonormal frame, orthonormalised each step), which explains why the columns settle onto the dominant eigenvectors. For symmetric $A$ the limit is diagonal.

\begin{nuance}{the practical algorithm}
The basic version is slow and costs $O(n^3)$ per step. Practical implementations (i) first reduce $A$ to \emph{Hessenberg form} (zeros below the first subdiagonal; tridiagonal if $A$ is symmetric) with $O(n^3)$ Householder reflections, so that each QR step costs only $O(n^2)$ (or $O(n)$ for tridiagonal); (ii) use \emph{shifts} $A_k-\mu_kI=Q_kR_k$, $A_{k+1}=R_kQ_k+\mu_kI$, which give quadratic or cubic convergence; and (iii) \emph{deflate}, splitting off converged eigenvalues. The total cost is about $\frac{10}3n^3$ flops for the eigenvalues of a symmetric matrix and $10n^3$ for a general one.
\end{nuance}

\section{Krylov subspaces}

Both power iteration and the QR algorithm use only products with $A$. The most powerful iterative methods use \emph{all} the information accumulated along the way.

\begin{definition}[Krylov subspace]
For $A\in\R^{n\times n}$ and $\mathbf{b}\in\R^n$, the \emph{Krylov subspace} is $\mathcal{K}_k(A,\mathbf{b})=\Span\{\mathbf{b},A\mathbf{b},A^2\mathbf{b},\dots,A^{k-1}\mathbf{b}\}$.
\end{definition}

Equivalently $\mathcal{K}_k=\{p(A)\mathbf{b}:\deg p<k\}$. By Cayley--Hamilton (Theorem~\ref{thm:ch}) $A^{-1}\mathbf{b}\in\mathcal{K}_n$, so the exact solution of $A\mathbf{x}=\mathbf{b}$ lies in $\mathcal{K}_n$; the hope is that a good approximation already lies in a much smaller $\mathcal{K}_k$.

\subsection*{Arnoldi and Lanczos}

The vectors $\mathbf{b},A\mathbf{b},\dots$ become nearly parallel, so we orthonormalise them with Gram--Schmidt (Chapter~11) as we go, obtaining an orthonormal basis $Q_k=[\mathbf{q}_1\cdots\mathbf{q}_k]$ of $\mathcal{K}_k$. This is the \emph{Arnoldi process}; it produces
\[ AQ_k=Q_{k+1}\tilde H_k,\qquad Q_k\T AQ_k=H_k, \]
with $H_k$ an upper Hessenberg $k\times k$ matrix, the compression of $A$ onto the Krylov subspace. When $A$ is symmetric, $H_k=T_k$ is tridiagonal (the \emph{Lanczos process}), so each new vector needs only a three-term recurrence $\beta_k\mathbf{q}_{k+1}=A\mathbf{q}_k-\alpha_k\mathbf{q}_k-\beta_{k-1}\mathbf{q}_{k-1}$. The eigenvalues of $H_k$ (\emph{Ritz values}) approximate the extreme eigenvalues of $A$ remarkably quickly, which is how large sparse eigenproblems are solved in practice (ARPACK, \texttt{scipy.sparse.linalg.eigsh}).

\section{The conjugate gradient method}

For a symmetric positive definite system $A\mathbf{x}=\mathbf{b}$, the best use of the Krylov subspace is the \emph{conjugate gradient} (CG) method of Hestenes and Stiefel. It minimises the quadratic $\phi(\mathbf{x})=\tfrac12\mathbf{x}\T A\mathbf{x}-\mathbf{b}\T\mathbf{x}$, whose gradient is $A\mathbf{x}-\mathbf{b}$ (Example~\ref{ex:lscalc}), by searching along $A$-orthogonal directions.

\begin{definition}[$A$-inner product, conjugacy]
For $A\succ0$ define $\ip{\mathbf{x}}{\mathbf{y}}_A=\mathbf{x}\T A\mathbf{y}$ and $\norm{\mathbf{x}}_A=\sqrt{\mathbf{x}\T A\mathbf{x}}$ (this is an inner product by Chapter~16). Vectors with $\mathbf{p}_i\T A\mathbf{p}_j=0$ are called \emph{conjugate} (A-orthogonal).
\end{definition}

\begin{theorem}[Conjugate gradient]\label{thm:cg}
Starting from $\mathbf{x}_0=\mathbf{0}$, $\mathbf{r}_0=\mathbf{p}_0=\mathbf{b}$, iterate for $k=0,1,\dots$
\[ \alpha_k=\frac{\mathbf{r}_k\T\mathbf{r}_k}{\mathbf{p}_k\T A\mathbf{p}_k},\quad\mathbf{x}_{k+1}=\mathbf{x}_k+\alpha_k\mathbf{p}_k,\quad\mathbf{r}_{k+1}=\mathbf{r}_k-\alpha_kA\mathbf{p}_k,\quad\beta_k=\frac{\mathbf{r}_{k+1}\T\mathbf{r}_{k+1}}{\mathbf{r}_k\T\mathbf{r}_k},\quad\mathbf{p}_{k+1}=\mathbf{r}_{k+1}+\beta_k\mathbf{p}_k. \]
Then the residuals $\mathbf{r}_j$ are mutually orthogonal, the directions $\mathbf{p}_j$ are mutually conjugate, $\mathbf{x}_k$ minimises $\norm{\mathbf{x}-\mathbf{x}^*}_A$ over $\mathcal{K}_k(A,\mathbf{b})$, and $\mathbf{x}_n=\mathbf{x}^*=A^{-1}\mathbf{b}$ in exact arithmetic (at most $n$ steps).
\end{theorem}

\begin{proof}[Idea]
One proves by induction that $\mathcal{K}_k=\Span\{\mathbf{r}_0,\dots,\mathbf{r}_{k-1}\}=\Span\{\mathbf{p}_0,\dots,\mathbf{p}_{k-1}\}$, that $\mathbf{r}_k\perp\mathcal{K}_k$ and that $\mathbf{p}_k$ is $A$-orthogonal to $\mathbf{p}_0,\dots,\mathbf{p}_{k-1}$; the recurrence for $\beta_k$ achieves this using only the most recent direction because of the symmetry of $A$. Residual orthogonality $\mathbf{r}_k\perp\mathcal{K}_k$ is exactly the condition that $\mathbf{x}_k$ is the $A$-orthogonal projection of $\mathbf{x}^*$ onto $\mathcal{K}_k$, hence the minimiser of the $A$-norm error (Theorem~\ref{thm:bestapprox} in the $A$-inner product). Orthogonal residuals in $\R^n$ can number at most $n$, so the residual vanishes after at most $n$ steps.
\end{proof}

\begin{theorem}[Convergence rate of CG]\label{thm:cgrate}
With $\kappa=\lambda_{\max}(A)/\lambda_{\min}(A)$,
\[ \norm{\mathbf{x}_k-\mathbf{x}^*}_A\le2\left(\frac{\sqrt\kappa-1}{\sqrt\kappa+1}\right)^k\norm{\mathbf{x}_0-\mathbf{x}^*}_A. \]
\end{theorem}

\begin{proof}[Idea]
$\mathbf{x}_k-\mathbf{x}^*=p_k(A)(\mathbf{x}_0-\mathbf{x}^*)$ for the polynomial $p_k$ of degree $\le k$ with $p_k(0)=1$ minimising the error in the $A$-norm. Using the spectral theorem, the error is at most $\max_{\lambda\in[\lambda_{\min},\lambda_{\max}]}\abs{p(\lambda)}$ for any such $p$; choosing $p$ to be a scaled Chebyshev polynomial gives the bound.
\end{proof}

\begin{interview}{CG versus gradient descent}
Gradient descent contracts the error by $\frac{\kappa-1}{\kappa+1}$ per step (Chapter~19); CG by $\frac{\sqrt\kappa-1}{\sqrt\kappa+1}$. For $\kappa=10^4$ that is the difference between $\approx10^4$ and $\approx10^2$ iterations. CG needs only one matrix--vector product and $O(n)$ extra storage per step, which is why it is the standard solver for large sparse SPD systems (finite elements, Gaussian processes, Newton-CG). Fast convergence needs \emph{preconditioning}: solve $M^{-1}A\mathbf{x}=M^{-1}\mathbf{b}$ with $M\approx A$ cheap to invert, so that $\kappa(M^{-1}A)\ll\kappa(A)$.
\end{interview}

For non-symmetric systems the analogues are GMRES (minimise the 2-norm of the residual over $\mathcal{K}_k$, using Arnoldi) and BiCGSTAB.

\section{Randomised SVD}

A full SVD costs $O(mn\min(m,n))$. When $A\in\R^{m\times n}$ is large and approximately low-rank, a randomised algorithm finds a good rank-$k$ approximation in $O(mnk)$ time.

\begin{theorem}[Randomised range finder, informal]\label{thm:randsvd}
Draw a Gaussian test matrix $\Omega\in\R^{n\times(k+p)}$ with small oversampling $p\approx5$--$10$. Form $Y=A\Omega$, an orthonormal basis $Q$ of $\Col Y$ (by QR), and $B=Q\T A\in\R^{(k+p)\times n}$. Compute the SVD $B=\tilde U\Sigma V\T$ and set $U=Q\tilde U$. Then $A\approx U\Sigma V\T$ and, in expectation, $\norm{A-QQ\T A}_2\le\big(1+O(\sqrt{k/p})\big)\sigma_{k+1}(A)$, nearly the optimum of Eckart--Young.
\end{theorem}

\begin{intuition}
A random combination of the columns of $A$ is very likely to have a component along each of the top singular directions, so $k+p$ random combinations span nearly the dominant column space; projecting onto it and taking a small SVD finishes the job. Power iterations $Y=(AA\T)^qA\Omega$ sharpen the spectrum when the singular values decay slowly.
\end{intuition}

\section{Worked interview problems}

The following problems are classical. Each is solved by a short argument using results from this book.

\begin{example}[Problem 1: rank of $A\T A$]
\emph{Show $\rank(A\T A)=\rank A$ for real $A$.} If $A\T A\mathbf{x}=\mathbf{0}$ then $\norm{A\mathbf{x}}^2=\mathbf{x}\T A\T A\mathbf{x}=0$, so $A\mathbf{x}=\mathbf{0}$. Hence $\Null A\T A=\Null A$ and rank--nullity gives equal ranks. (This fails over $\C$ for $A\T A$ with transpose instead of conjugate transpose, and over finite fields.)
\end{example}

\begin{example}[Problem 2: eigenvalues of a rank-one matrix]
\emph{Find the eigenvalues of $\mathbf{u}\mathbf{v}\T$.} The matrix has rank at most $1$, so $0$ has geometric multiplicity at least $n-1$. The nonzero eigenvalue, if any, satisfies $\tr(\mathbf{u}\mathbf{v}\T)=\mathbf{v}\T\mathbf{u}$. So the eigenvalues are $\mathbf{v}\T\mathbf{u}$ (once) and $0$ (with multiplicity $n-1$). Check: $(\mathbf{u}\mathbf{v}\T)\mathbf{u}=(\mathbf{v}\T\mathbf{u})\mathbf{u}$. It is diagonalisable iff $\mathbf{v}\T\mathbf{u}\ne0$ (otherwise it is nilpotent and nonzero).
\end{example}

\begin{example}[Problem 3: $AB$ versus $BA$]
\emph{Show $AB$ and $BA$ have the same characteristic polynomial for square $A,B$ (and the same nonzero eigenvalues if rectangular).} If $\lambda\ne0$ and $AB\mathbf{v}=\lambda\mathbf{v}$ with $\mathbf{v}\ne\mathbf{0}$, then $B\mathbf{v}\ne\mathbf{0}$ and $BA(B\mathbf{v})=B(AB\mathbf{v})=\lambda B\mathbf{v}$, so $\lambda$ is an eigenvalue of $BA$. For square matrices, let $M=\begin{pmatrix}\lambda I&A\\B&I\end{pmatrix}$ with $\lambda\ne0$. Taking the Schur complement of the lower-right block $I$ gives $\det M=\det(\lambda I-AB)$, and taking the Schur complement of the upper-left block $\lambda I$ gives $\det M=\lambda^n\det(I-BA/\lambda)=\det(\lambda I-BA)$. So $p_{AB}(\lambda)=p_{BA}(\lambda)$ for all $\lambda\ne0$, hence as polynomials.
\end{example}

\begin{example}[Problem 4: nilpotent matrices]
\emph{If $N^k=0$, prove $N$ has only the eigenvalue $0$, $I+N$ is invertible, and $N^n=0$.} If $N\mathbf{v}=\lambda\mathbf{v}$ then $0=N^k\mathbf{v}=\lambda^k\mathbf{v}$, so $\lambda=0$. Then $p_N=\lambda^n$ and Cayley--Hamilton gives $N^n=0$. Finally $(I+N)^{-1}=I-N+N^2-\dots\pm N^{n-1}$ (Neumann series, finite).
\end{example}

\begin{example}[Problem 5: trace and eigenvalues of a projection]
\emph{If $P^2=P$, show $\rank P=\tr P$.} The eigenvalues of $P$ are $0$ or $1$ (from $\lambda^2=\lambda$), and $P$ is diagonalisable since its minimal polynomial $\lambda^2-\lambda$ has distinct roots (Theorem~\ref{thm:minpoly}). So $\tr P=$ number of eigenvalues equal to $1=\dim\Col P=\rank P$.
\end{example}

\begin{example}[Problem 6: positive definite from a sum]
\emph{If $A\succ0$ and $B\succeq0$, show $A+B\succ0$ and $\det(A+B)\ge\det A$.} $\mathbf{x}\T(A+B)\mathbf{x}\ge\mathbf{x}\T A\mathbf{x}>0$. For the determinant, $\det(A+B)=\det A\,\det(I+A^{-1/2}BA^{-1/2})$ and the second factor is the product of $1+\mu_i$ with $\mu_i\ge0$ the eigenvalues of the PSD matrix $A^{-1/2}BA^{-1/2}$, hence $\ge1$.
\end{example}

\begin{example}[Problem 7: the Cayley transform]
\emph{If $S$ is skew-symmetric, prove $(I-S)^{-1}(I+S)$ is orthogonal.} $I-S$ is invertible since $\mathbf{x}\T(I-S)\mathbf{x}=\norm{\mathbf{x}}^2>0$. Let $Q=(I-S)^{-1}(I+S)$. Using $S\T=-S$ and that $(I\pm S)$ commute, $Q\T=(I-S)(I+S)^{-1}$ so $Q\T Q=(I-S)(I+S)^{-1}(I-S)^{-1}(I+S)=I$.
\end{example}

\begin{example}[Problem 8: sum of squares of eigenvalues]
\emph{For symmetric $A$, show $\sum\lambda_i^2=\sum_{i,j}a_{ij}^2$.} $\tr(A^2)=\sum\lambda_i^2$ (eigenvalues of $A^2$) and $\tr(A\T A)=\sum_{ij}a_{ij}^2$; for symmetric $A$ these coincide, i.e.\ $\norm{A}_F^2=\sum\lambda_i^2$.
\end{example}

\begin{example}[Problem 9: matrix with prescribed row and column space]
\emph{Is there a $3\times3$ matrix with $\Col A=\Span\{(1,0,0),(0,1,0)\}$ and $\Null A=\Span\{(1,1,1)\}$?} Rank--nullity forces $\rank A=2$ and $\dim\Null A=1$, consistent. Build it: choose $A\mathbf{e}_1=\mathbf{e}_1$, $A\mathbf{e}_2=\mathbf{e}_2$ and $A(1,1,1)\T=\mathbf{0}$, so $A\mathbf{e}_3=-\mathbf{e}_1-\mathbf{e}_2$. Then $A=\begin{pmatrix}1&0&-1\\0&1&-1\\0&0&0\end{pmatrix}$.
\end{example}

\begin{example}[Problem 10: distance to a subspace]
\emph{Find the distance from $\mathbf{b}=(1,2,3,4)$ to $\Span\{(1,1,0,0),(0,0,1,1)\}$.} The projection is $\mathbf{p}=\frac32(1,1,0,0)+\frac72(0,0,1,1)=(1.5,1.5,3.5,3.5)$ (the subspace has orthogonal basis $\mathbf{a}_1,\mathbf{a}_2$), so the residual is $(-0.5,0.5,-0.5,0.5)$ with norm $1$.
\end{example}

\begin{example}[Problem 11: powers of a stochastic-like matrix]
\emph{Compute $\lim_{k\to\infty}A^k$ for $A=\begin{pmatrix}0.5&0.5\\0.25&0.75\end{pmatrix}$.} $A$ is stochastic and regular with $\boldsymbol\pi A=\boldsymbol\pi$: $0.5\pi_1+0.25\pi_2=\pi_1$ gives $\pi_2=2\pi_1$, so $\boldsymbol\pi=(\tfrac13,\tfrac23)$ and $A^k\to\begin{pmatrix}1/3&2/3\\1/3&2/3\end{pmatrix}$.
\end{example}

\begin{example}[Problem 12: a determinant identity]
\emph{Show $\det(I-\mathbf{u}\mathbf{v}\T)=1-\mathbf{v}\T\mathbf{u}$.} The eigenvalues of $\mathbf{u}\mathbf{v}\T$ are $\mathbf{v}\T\mathbf{u}$ and zeros (Problem 2), so those of $I-\mathbf{u}\mathbf{v}\T$ are $1-\mathbf{v}\T\mathbf{u}$ and ones, and the determinant is their product.
\end{example}

\begin{pitfall}{common interview traps}
(1) Claiming $\det(A+B)=\det A+\det B$ or $(AB)^{-1}=A^{-1}B^{-1}$. (2) Forgetting that eigenvectors must be nonzero and that a defective matrix has fewer eigenvectors than its size. (3) Saying ``symmetric implies distinct eigenvalues''; it implies \emph{real} eigenvalues and an orthonormal eigenbasis. (4) Computing the inverse to solve a system. (5) Confusing the rank of a matrix with the number of nonzero \emph{eigenvalues} for non-symmetric matrices ($\begin{pmatrix}0&1\\0&0\end{pmatrix}$ has rank $1$ and no nonzero eigenvalues). (6) Thinking that positive determinant and positive diagonal imply positive definiteness. (7) Forgetting to centre data before PCA.
\end{pitfall}

\section{A map of the subject}

\begin{center}
\begin{tabular}{p{3.6cm}p{9cm}}
\toprule
Question & Tool \\
\midrule
Does $A\mathbf{x}=\mathbf{b}$ have a solution? How many? & Pivots, rank, $\Col A$, $\Null A$ (Chapters 2, 8) \\
Solve repeatedly or accurately & $PA=LU$, Cholesky, QR (Chapters 4, 11) \\
Best fit for an overdetermined system & Projection, normal equations, QR, SVD (Chapters 12, 17) \\
What does a matrix do geometrically? & Determinant, eigenvectors, SVD (Chapters 5, 13, 17) \\
Powers, exponentials, stability & Diagonalisation, Jordan, spectral radius (Chapters 14, 18) \\
Symmetric structure, optimisation & Spectral theorem, PSD, Rayleigh (Chapters 15, 16) \\
Compress or denoise & Truncated SVD, Eckart--Young (Chapter 17) \\
Numerical trust & $\kappa(A)$, backward stability (Chapter 18) \\
Random walks and networks & Perron--Frobenius, Laplacians (Chapter 20) \\
Large sparse problems & Power, Lanczos, CG (this chapter) \\
\bottomrule
\end{tabular}
\end{center}

\begin{summarybox}
\begin{itemize}
  \item Power iteration finds the dominant eigenpair using only matrix--vector products; inverse iteration targets eigenvalues near a shift, and Rayleigh quotient iteration converges cubically for symmetric matrices.
  \item The QR algorithm iterates $A_k=Q_kR_k$, $A_{k+1}=R_kQ_k$; with Hessenberg reduction and shifts it computes all eigenvalues in $O(n^3)$.
  \item Krylov subspaces $\Span\{\mathbf{b},A\mathbf{b},\dots\}$ underlie Arnoldi, Lanczos and GMRES; CG minimises the $A$-norm error over them, converges at rate $\frac{\sqrt\kappa-1}{\sqrt\kappa+1}$ and finishes in at most $n$ steps.
  \item Randomised range finders give near-optimal low-rank approximations in $O(mnk)$ time.
  \item The classic interview problems reduce to rank--nullity, trace and determinant identities, eigenvalue arguments, projections and positive definiteness.
\end{itemize}
\end{summarybox}

\section{Exercises}

\begin{exercise}
Apply three steps of power iteration to $A=\begin{pmatrix}4&1\\2&3\end{pmatrix}$ starting at $(1,1)$, and compare with the exact dominant eigenvalue $5$ and eigenvector.
\end{exercise}

\begin{exercise}
Perform two steps of the basic QR algorithm on $A=\begin{pmatrix}2&1\\1&3\end{pmatrix}$ and observe the off-diagonal entry shrinking.
\end{exercise}

\begin{exercise}
Run conjugate gradient by hand on $A=\begin{pmatrix}4&1\\1&3\end{pmatrix}$, $\mathbf{b}=(1,2)\T$, and verify that it terminates in two steps with $\mathbf{x}=(1/11,7/11)$.
\end{exercise}

\begin{exercise}
Show that if $\abs{\lambda_1}=\abs{\lambda_2}$ with $\lambda_1\ne\lambda_2$, the power method need not converge, and give an example.
\end{exercise}

\begin{exercise}
Prove that for symmetric tridiagonal $T_k$ produced by Lanczos, the eigenvalues of $T_k$ interlace with those of $T_{k+1}$ (Cauchy interlacing).
\end{exercise}

\begin{exercise}
Estimate the number of CG iterations needed to reduce the $A$-norm error by $10^{-6}$ when $\kappa=10^4$, and compare with gradient descent.
\end{exercise}

\begin{exercise}
Prove Sylvester's determinant theorem $\det(I+AB)=\det(I+BA)$ for $A\in\R^{m\times n}$, $B\in\R^{n\times m}$, using the block factorisations of $\begin{pmatrix}I&-A\\B&I\end{pmatrix}$, and deduce Problem 12.
\end{exercise}

\begin{exercise}
Design your own interview question that uses (a) rank--nullity, (b) the spectral theorem and (c) a trace identity, and write its solution.
\end{exercise}
